%% composable-future.tex
%% Composable Future: Toward an Algebraic Theory of Paradigmatic Transitions
%% Target: arXiv cs.LO / math.CT
%% Compile: pdflatex → biber → pdflatex × 2
%%          or: latexmk -pdf composable-future.tex

\documentclass[11pt, a4paper]{article}

% ── Encoding & fonts ──────────────────────────────────────────────────────────
\usepackage[T1]{fontenc}
\usepackage[utf8]{inputenc}
\usepackage{lmodern}
\usepackage{microtype}

% ── Page geometry ─────────────────────────────────────────────────────────────
\usepackage[
  top=2.5cm, bottom=2.5cm,
  left=3cm,  right=3cm
]{geometry}

% ── Mathematics ───────────────────────────────────────────────────────────────
\usepackage{amsmath}
\usepackage{amssymb}
\usepackage{amsthm}
\usepackage{mathtools}   % for \coloneqq, prescript, etc.
\usepackage{stmaryrd}    % for \llbracket, \rrbracket

% ── Diagrams (optional — comment out if not compiling locally) ────────────────
% \usepackage{tikz-cd}   % for commutative diagrams

% ── References ────────────────────────────────────────────────────────────────
\usepackage[
  backend=biber,
  style=authoryear,
  maxcitenames=2,
  maxbibnames=99,
  doi=true,
  url=false
]{biblatex}
\addbibresource{composable-future.bib}

% ── Hyperlinks ────────────────────────────────────────────────────────────────
\usepackage[
  colorlinks=true,
  linkcolor=black,
  citecolor=black,
  urlcolor=black
]{hyperref}

% ── Miscellaneous ─────────────────────────────────────────────────────────────
\usepackage{enumitem}
\usepackage{booktabs}    % for \toprule etc. in tables
\usepackage{xcolor}
\usepackage{setspace}

% ── Theorem environments ──────────────────────────────────────────────────────
\theoremstyle{plain}
\newtheorem{theorem}{Theorem}[section]
\newtheorem{proposition}[theorem]{Proposition}
\newtheorem{lemma}[theorem]{Lemma}
\newtheorem{corollary}[theorem]{Corollary}

\theoremstyle{definition}
\newtheorem{definition}[theorem]{Definition}
\newtheorem{example}[theorem]{Example}

\theoremstyle{remark}
\newtheorem{remark}[theorem]{Remark}
\newtheorem{openproblem}{Open Problem}

% ── Notation ──────────────────────────────────────────────────────────────────

% Operators
\newcommand{\bind}{\mathbin{\gg\!=}}          % sequential: >>=
\newcommand{\tensor}{\otimes}                  % parallel:   ⊗
\newcommand{\pfork}{\mathbin{|}}               % fork:       |
\newcommand{\merge}{\oplus}                    % merge:      ⊕
\newcommand{\nullfut}{\mathrm{Id}}             % null future

% Sets and types
\newcommand{\Fut}{\mathcal{F}}                 % set of all futures
\newcommand{\Pow}[1]{\mathcal{P}(#1)}         % power set
\newcommand{\States}{\mathcal{S}}              % set of paradigmatic states
\newcommand{\Aff}[1]{\Phi_{#1}}               % affordance set of F

% Components of a future
\newcommand{\Szero}{S_0}                       % initial state
\newcommand{\Sone}{S_1}                        % realized state
\newcommand{\traj}{\tau}                       % trajectory morphism
\newcommand{\aff}{\Phi}                        % affordance set

% Category-theoretic
\newcommand{\Cat}{\mathbf{Cat}}
\newcommand{\Set}{\mathbf{Set}}
\newcommand{\Stoch}{\mathbf{Stoch}}            % Markov/stochastic category

% Convenience
\newcommand{\ie}{i.e.,\xspace}
\newcommand{\eg}{e.g.,\xspace}
\newcommand{\cf}{cf.\xspace}

% ── Title block ───────────────────────────────────────────────────────────────
\title{%
  \textbf{Composable Future}\\[4pt]
  \large Toward an Algebraic Theory of Paradigmatic Transitions
}

\author{%
  \textit{[Author]}\\[2pt]
  \small \textit{[Institution]}\\
  \small \href{mailto:}{[email]}
}

\date{%
  \small Version 0.1 — \today\\
  \small Preprint. Comments welcome.
}

% ─────────────────────────────────────────────────────────────────────────────
\begin{document}
% ─────────────────────────────────────────────────────────────────────────────

\maketitle

\begin{abstract}
We introduce \emph{Composable Future}, a formal structure in which paradigmatic
futures are treated as algebraic objects admitting composition.
A future $F$ is defined as a 4-tuple $(S_0, \traj, S_1, \aff)$ where $S_0$
and $S_1$ are paradigmatic states, $\traj\colon S_0 \to S_1$ is a trajectory
morphism, and $\aff\colon S_1 \to \Pow{\Fut}$ is a typed affordance set.
We define four primitive operators---sequential bind $(\bind)$, parallel
tensor $(\tensor)$, fork $(\pfork)$, and merge $(\merge)$---and investigate
the algebraic laws they satisfy.
We show that identity and closure hold in general; associativity of $\bind$
holds when $\traj$ is stateless, and reduces to a fibered structure when
$\traj$ is path-dependent.
Under probabilistic extension, the structure forms a Kleisli category over
the probability monad.
We identify five open problems and position the framework relative to applied
category theory, process algebra, branching-time temporal logic, and
affordance theory.
\end{abstract}

\smallskip
\noindent\textbf{Keywords:}
composable future, paradigmatic transition, category theory, process algebra,
affordance, indexed monad, branching time.

\bigskip

% ─────────────────────────────────────────────────────────────────────────────
\section{Introduction}
\label{sec:intro}
% ─────────────────────────────────────────────────────────────────────────────

\subsection{The problem}
% WRITE: Paradigmatic transitions are currently modeled only descriptively.
% Kuhn (1962) identified them. Lakatos (1978) extended the taxonomy.
% Neither produced a formal algebra.
% The question of how paradigms compose has not been posed formally.
%
% Key move: distinguish "paradigm shift" (event-level) from
% "paradigmatic transition" (structural-level) — we are formalizing the latter.

\textbf{[TODO: \S1.1 — motivate the gap; 2--3 paragraphs]}

\subsection{What is missing}
% WRITE: Three adjacent formalisms come close:
% (1) Applied CT (Spivak 2014, Furter et al. 2025) — within systems, not between paradigmatic states
% (2) Process algebra (Hoare 1985, Katis et al. 2009) — concurrent processes, not paradigmatic transitions
% (3) Branching-time logic (Iacona & Iaquinto 2021) — truth over futures, not algebraic composition
%
% Table format works well here.

\textbf{[TODO: \S1.2 — three adjacent formalisms, each with one sentence on
what they cover and one on what they miss; table optional]}

\subsection{Contribution}
% WRITE: What this paper does and does not do.
% Does: defines F, four operators, states laws, identifies open problems.
% Does not: prove associativity (open problem P1), mechanize (future work).

\textbf{[TODO: \S1.3 — precise statement of contributions; list format]}

% ─────────────────────────────────────────────────────────────────────────────
\section{The Primitive Object}
\label{sec:primitive}
% ─────────────────────────────────────────────────────────────────────────────

\subsection{Paradigmatic states}

\begin{definition}[Paradigmatic state]
\label{def:state}
A \emph{paradigmatic state} $S \in \States$ is a structured triple
$(A, C, I)$ where:
\begin{itemize}[itemsep=2pt]
  \item $A$ is a set of foundational assumptions operative in the domain;
  \item $C$ is a set of constraints on valid operations and inferences;
  \item $I$ is a set of infrastructure elements enabling those operations.
\end{itemize}
\end{definition}

\begin{remark}
Definition~\ref{def:state} is intentionally domain-agnostic.
The triple $(A, C, I)$ is not required to be convex in any geometric sense
(\cf\ the conceptual space formalization of
\textcite{bechberger2017thorough}), as paradigmatic states need not occupy
contiguous regions in any similarity space.
\end{remark}

\textbf{[TODO: 1 paragraph — note that S is an object in the category we
are building; forward reference to \S4]}

\subsection{The Composable Future}

\begin{definition}[Composable Future]
\label{def:future}
A \emph{Composable Future} is a 4-tuple
\[
  F = (\Szero,\; \traj,\; \Sone,\; \aff)
\]
where:
\begin{enumerate}[label=(\roman*), itemsep=3pt]
  \item $\Szero \in \States$ is the \emph{initial paradigmatic state};
  \item $\traj\colon \Szero \to \Sone$ is the \emph{trajectory}, a morphism
        encoding the mechanism of transition;
  \item $\Sone \in \States$ is the \emph{reachable paradigmatic state};
  \item $\aff\colon \Sone \to \Pow{\Fut}$ is the \emph{affordance set},
        a function from the realized state to the set of futures it makes
        compositionally accessible.
\end{enumerate}
We denote the collection of all Composable Futures by $\Fut$.
\end{definition}

\subsection{The null future}

\begin{definition}[Null future]
\label{def:null}
For any paradigmatic state $S \in \States$, the \emph{null future} at $S$
is the Composable Future
\[
  \nullfut_S \;\coloneqq\; (S,\; \mathrm{id}_S,\; S,\; \aff_\varnothing)
\]
where $\mathrm{id}_S$ is the identity morphism on $S$ and
$\aff_\varnothing(S) = \varnothing$ for all $S$.
\end{definition}

\textbf{[TODO: 1 paragraph — role of null future in stating identity law;
forward reference to Proposition~\ref{prop:identity}]}

\subsection{Affordances}

% WRITE: Ground Φ philosophically in Chemero (2003) relational ontology.
% Φ is not a property of S₁ alone — it is a relation between S₁ and the
% composition context.
% Reference Turvey (1992) for prospective / forward-directed nature.
% Note that Φ is not well-defined before S₁ is realized — motivates §6.

\textbf{[TODO: \S2.4 — philosophical grounding of $\aff$; 2 paragraphs;
cite Chemero 2003, Turvey 1992]}

% ─────────────────────────────────────────────────────────────────────────────
\section{Operators}
\label{sec:operators}
% ─────────────────────────────────────────────────────────────────────────────

We define four primitive operations over $\Fut$.
Each operator takes one or two Composable Futures and returns a Composable
Future, subject to compatibility conditions stated below.

\subsection{Sequential composition}

\begin{definition}[Sequential bind]
\label{def:bind}
Let $A = (\Szero^A, \traj^A, \Sone^A, \aff^A)$ and
$B = (\Szero^B, \traj^B, \Sone^B, \aff^B)$ be Composable Futures
with $\Sone^A = \Szero^B$.
Their \emph{sequential composition} is
\[
  A \bind B \;\coloneqq\;
  \bigl(\Szero^A,\; \traj^A ; \traj^B,\; \Sone^B,\; \aff^B\bigr)
\]
where $\traj^A ; \traj^B$ denotes composition of morphisms in $\States$.
\end{definition}

\begin{remark}
The affordance set of $A$ is consumed by the composition: only $\aff^B$
appears in the result. Intuitively, the affordances of $A$ are realized
through $B$ rather than carried forward independently.
\end{remark}

\textbf{[TODO: 1 sentence — note that the compatibility condition
$\Sone^A = \Szero^B$ parallels interface matching in TCP
\parencite{katis2009process}]}

\subsection{Parallel tensor}

\begin{definition}[Parallel tensor]
\label{def:tensor}
Let $A$ and $B$ be Composable Futures.
Their \emph{parallel tensor} is
\[
  A \tensor B \;\coloneqq\;
  \bigl(\Szero^A \otimes \Szero^B,\;
        \traj^A \otimes \traj^B,\;
        \Sone^A \otimes \Sone^B,\;
        \aff^A \times \aff^B\bigr)
\]
where $\otimes$ on states is the monoidal product on $\States$ and
$\aff^A \times \aff^B$ denotes the product affordance set.
\end{definition}

% WRITE: Note grounding in Span(Graph) SMC structure (Katis et al. 2009)
% and co-design tensor (Furter et al. 2025).

\textbf{[TODO: 1--2 sentences grounding in Span(Graph) and SMC literature]}

\subsection{Fork}

\begin{definition}[Fork]
\label{def:fork}
Let $A$ and $B$ be Composable Futures with $\Szero^A = \Szero^B = S$.
Their \emph{fork} is
\[
  A \pfork B \;\coloneqq\;
  \bigl(S,\; \traj^A \sqcup \traj^B,\; \Sone^A \sqcup \Sone^B,\;
        \aff^A \sqcup \aff^B\bigr)
\]
where $\sqcup$ denotes the coproduct (disjoint union) on states and
$\traj^A \sqcup \traj^B$ is the induced morphism out of $S$.
Exactly one branch is realized at any history.
\end{definition}

% WRITE: Corresponds to nondeterministic choice □ in CSP (Hoare 1985)
% and branching-time fork of Iacona & Iaquinto (2021).

\textbf{[TODO: 1 sentence connecting to CSP $\square$ and branching-time]}

\subsection{Merge}

\begin{definition}[Merge]
\label{def:merge}
Let $A$ and $B$ be Composable Futures with compatible realized states,
\ie\ $\Sone^A \sim \Sone^B$ under a compatibility relation $\sim$ to be
specified.
Their \emph{merge} is
\[
  A \merge B \;\coloneqq\;
  \bigl(\Szero^A \times \Szero^B,\;
        \traj^A \times \traj^B,\;
        \Sone^A \sqcap \Sone^B,\;
        \aff^A \cap \aff^B\bigr)
\]
where $\sqcap$ is the meet on compatible states.
\end{definition}

\begin{remark}
The compatibility relation $\sim$ is left as a parameter.
A natural choice is bisimulation \parencite{milner1989communication}:
$\Sone^A \sim \Sone^B$ iff they behave identically under all contexts.
\end{remark}

% ─────────────────────────────────────────────────────────────────────────────
\section{Algebraic Laws}
\label{sec:laws}
% ─────────────────────────────────────────────────────────────────────────────

\subsection{Identity}

\begin{proposition}[Identity law]
\label{prop:identity}
For any Composable Future $F = (\Szero, \traj, \Sone, \aff)$:
\[
  F \bind \nullfut_{\Sone} \;=\; F
  \qquad\text{and}\qquad
  \nullfut_{\Szero} \bind F \;=\; F.
\]
\end{proposition}

\begin{proof}
Direct from Definition~\ref{def:bind}.
For the right identity:
$F \bind \nullfut_{\Sone}
  = (\Szero,\; \traj ; \mathrm{id}_{\Sone},\; \Sone,\; \aff_\varnothing)
  = (\Szero,\; \traj,\; \Sone,\; \aff_\varnothing)$.
The affordance set reduces because $\aff_\varnothing(\Sone) = \varnothing$;
the trajectory is preserved by the unit law of morphism composition.
Left identity is symmetric.
\qed
\end{proof}

\begin{remark}
The identity law confirms that $\nullfut_S$ plays the role of the identity
morphism in any category of Composable Futures.
Note that the affordance set of the result is $\aff_\varnothing$, not $\aff$.
This is expected: composing with the null future collapses the affordance
set, since no new futures are made accessible by a transition that changes
nothing.
\end{remark}

\subsection{Closure}

\begin{proposition}[Closure under $\bind$]
\label{prop:closure}
If $A, B \in \Fut$ and $\Sone^A = \Szero^B$, then $A \bind B \in \Fut$.
\end{proposition}

\begin{proof}
Immediate from Definition~\ref{def:bind}: the result is a 4-tuple
$(\Szero^A, \traj^A ; \traj^B, \Sone^B, \aff^B)$.
Each component is well-defined: $\Szero^A, \Sone^B \in \States$;
$\traj^A ; \traj^B$ is a valid morphism by composition closure of $\States$;
$\aff^B \colon \Sone^B \to \Pow{\Fut}$ is well-typed.
\qed
\end{proof}

\textbf{[TODO: 1 sentence — note analogous closure for $\tensor$, $\pfork$,
$\merge$ holds under their respective compatibility conditions]}

\subsection{Associativity — the open problem}

\begin{openproblem}[Associativity of $\bind$]
\label{op:assoc}
Does the following hold for all compatible $A, B, C \in \Fut$?
\[
  (A \bind B) \bind C \;=\; A \bind (B \bind C)
\]
\end{openproblem}

We distinguish two cases.

\medskip\noindent\textbf{Case 1 — stateless trajectory.}
If $\traj$ is stateless (\ie\ $\traj^B$ does not depend on whether $\traj^A$
has been executed), then morphism composition in $\States$ is associative
and the identity holds.
In this case $(\Fut, \bind, \nullfut)$ forms a \emph{category}.

\medskip\noindent\textbf{Case 2 — path-dependent trajectory.}
If $\traj^B$ is shaped by the prior execution of $\traj^A$ --- as is plausible
when the mechanism of change in one paradigmatic transition alters the
conditions under which a subsequent transition occurs --- then the grouping
$(A \bind B) \bind C \neq A \bind (B \bind C)$ in general, because
$\traj^B$ carries different semantics depending on whether $\traj^A$ has
been indexed into the composition or not.

\medskip\noindent\textbf{Candidate resolution.}
\textcite{orchard2014semantic} show that path-dependent composition can be
handled via \emph{indexed monads}, in which the monad carries an index
tracking prior computational effects.
Lifting this construction to paradigmatic trajectories would replace
$\traj\colon \Szero \to \Sone$ with an indexed morphism
$\traj_h\colon \Szero \to \Sone$ where $h$ records the history of prior
transitions.
Associativity would then hold at the level of the indexed structure even
when individual morphisms are history-carrying.
Full development is deferred to subsequent work.

\subsection{Commutativity of $\tensor$}

\begin{proposition}[Non-commutativity of $\tensor$]
\label{prop:noncomm}
In general, $A \tensor B \not\cong B \tensor A$.
\end{proposition}

% WRITE: Proof sketch — the affordance set Φ^A × Φ^B differs from
% Φ^B × Φ^A because the order of paradigmatic development changes what
% becomes compositionally accessible from S₁.
% This is not a deficiency but a structural feature.

\textbf{[TODO: proof sketch — 3--4 sentences; note that non-commutativity
is meaningful, not a failure]}

\subsection{Structural targets}

\begin{table}[h]
\centering
\caption{Algebraic structural targets for $(\Fut, \bind, \nullfut)$.}
\label{tab:structures}
\smallskip
\begin{tabular}{llll}
\toprule
Structure & Conditions required & Holds? \\
\midrule
Category      & Identity + associativity + closure & Case 1 \\
Monoid        & Category + single object            & Under investigation \\
Monad         & Monoid + \texttt{return} + associativity & Requires Case 1 or indexed resolution \\
Fibered cat.  & Path-dependent $\traj$              & Case 2 \\
\bottomrule
\end{tabular}
\end{table}

% ─────────────────────────────────────────────────────────────────────────────
\section{The Affordance Set}
\label{sec:affordance}
% ─────────────────────────────────────────────────────────────────────────────

\subsection{Philosophical grounding}

% WRITE: Chemero (2003) — affordances as relations, not properties.
% Φ is a relation between S₁ and the composition context.
% Turvey (1992) — prospective, forward-directed.
% Şahin et al. (2007) — formal triple (entity, behavior, effect) at robotics level.
% Kyriakoullis & Zaphiris (2016) — typed predicates, domain-dependence.
% Key claim: Φ is paradigm-specific — its type depends on the domain
% instantiating the theory.

\textbf{[TODO: \S5.1 — 2 paragraphs; cite Chemero, Turvey, Şahin et al.]}

\subsection{Formal treatment}

In dependent type theory, the affordance set is typed as:
\[
  \aff(s \colon \Sone) \;\colon\; \mathsf{Type}
\]
— the type of futures available given the realized state $s$.
This is a \emph{dependent type} over $\Sone$, which captures the
paradigm-specificity of affordances: different realized states yield
structurally different sets of accessible futures.

\subsection{Composition of affordances}

\begin{openproblem}[Affordance composition]
\label{op:aff-comp}
Does the following hold?
\[
  (A \bind B).\aff \;=\; \aff^B \circ \aff^A
\]
That is, are the affordances made accessible by $A \bind B$ fully
determined by composing the affordance functions of $A$ and $B$?
\end{openproblem}

% WRITE: Note realizability constraint — Φ is not well-defined before S₁
% is realized. Motivate §6 (probabilistic extension).

\textbf{[TODO: 1 paragraph — realizability constraint; forward to \S6]}

% ─────────────────────────────────────────────────────────────────────────────
\section{Probabilistic Extension}
\label{sec:probabilistic}
% ─────────────────────────────────────────────────────────────────────────────

When $\Sone$ is not a point but a distribution over paradigmatic states,
the future becomes:
\[
  F \;=\; \bigl(\Szero,\; \traj,\; \mathcal{P}(\Sone),\; \aff\bigr)
\]
where $\mathcal{P}(\Sone)$ is a probability measure over $\Sone$ and
$\traj\colon \Szero \to \mathcal{P}(\Sone)$ is a Markov kernel.

This is precisely the change-of-base construction for enriched categories
developed by \textcite{furter2025composable}, here applied to paradigmatic
states rather than engineering design components.
The resulting structure is a \textbf{Kleisli category} over the probability
monad $\mathcal{P}$: morphisms are Markov kernels, and sequential composition
$\bind$ becomes monadic bind in the Kleisli sense.

% WRITE: 2 additional sentences — note this resolves the realizability
% issue for Φ (Φ is typed over the distribution, not a point), and
% note that full development is deferred.

\textbf{[TODO: 2 sentences on Φ over distributions; defer full development]}

% ─────────────────────────────────────────────────────────────────────────────
\section{Open Problems}
\label{sec:open}
% ─────────────────────────────────────────────────────────────────────────────

We state five open problems that constitute the research programme
following from this paper.

\begin{openproblem}[Associativity under path-dependence]
\label{op:1}
Does $(A \bind B) \bind C = A \bind (B \bind C)$ hold when $\traj$ is
path-dependent?
The indexed monad construction of \textcite{orchard2014semantic} is the
primary candidate for resolution.
\end{openproblem}

\begin{openproblem}[Realizability of $\aff$]
\label{op:2}
Is $\aff$ well-defined before $\Sone$ is realized?
Under the probabilistic extension (\S\ref{sec:probabilistic}), $\aff$ may
be typed over the distribution $\mathcal{P}(\Sone)$ rather than a point,
which may resolve this constraint.
\end{openproblem}

\begin{openproblem}[Equivalence relation on futures]
\label{op:3}
What is the correct notion of equivalence between two futures $F$ and $F'$?
The candidate is \emph{bisimulation} \parencite{milner1989communication}:
$F \sim F'$ iff they behave identically under all composition contexts,
regardless of internal structure.
\end{openproblem}

\begin{openproblem}[Affordance composition]
\label{op:4}
Does $(A \bind B).\aff = \aff^B \circ \aff^A$ hold in general?
Under what conditions on $\States$ and $\aff$ does affordance composition
preserve the affordance type?
\end{openproblem}

\begin{openproblem}[Completeness]
\label{op:5}
Are all paradigmatic futures reachable by finite composition from a set of
primitive futures?
Formally: does there exist a finite generating set
$G \subset \Fut$ such that every $F \in \Fut$ is expressible as a finite
application of $\bind$, $\tensor$, $\pfork$, $\merge$ over $G$?
\end{openproblem}

% ─────────────────────────────────────────────────────────────────────────────
\section{Conclusion}
\label{sec:conclusion}
% ─────────────────────────────────────────────────────────────────────────────

% WRITE:
% (1) Summary of what was done — defined F, four operators, proved identity
%     and closure, stated laws.
% (2) Structural result — category in stateless case, fibered/indexed in
%     path-dependent case.
% (3) Relationship to five existing formalisms — CT, process algebra,
%     temporal logic, affordance theory, conceptual spaces.
% (4) The five open problems as a research programme.
% (5) One sentence on the mechanized proof target (Lean 4 / Agda).

\textbf{[TODO: \S8 — 3--4 paragraphs; write after all other sections
are complete]}

% ─────────────────────────────────────────────────────────────────────────────
\printbibliography
% ─────────────────────────────────────────────────────────────────────────────

\end{document}
